\section{Project Specification}
\label{sec:specification}

\begin{guidance}[What this section does -- about 200 words plus the table]
A detailed product or process design specification, justified by the
literature, the market research or a user survey. Use short, definitive
statements and quantify every parameter you can. Number the requirements
with \verb|\req{req:label}|: the IDs R1, R2 \ldots\ can then be used as
evaluation criteria in the options analysis (\cref{sec:options}) and, later,
in the testing and evaluation of the Final Report. See \textcite{cross2008}
and \textcite{ulrich2012}; on testable requirements \textcite{dick2017}.
\end{guidance}

\begin{table}[htbp]
  \centering
  \caption{Product design specification (example)}
  \label{tab:requirements}
  \small
  \begin{tabular}{p{0.06\linewidth}p{0.14\linewidth}p{0.36\linewidth}cp{0.2\linewidth}}
    \toprule
    ID & Category & Requirement & D/W & Source \\
    \midrule
    \req{req:accuracy} & Performance & Temperature measured with an accuracy
      of $\pm\SI{0.5}{\celsius}$ between \SIrange{0}{40}{\celsius}
      & D & Market survey (\cref{tab:market}) \\
    \req{req:interval} & Performance & Logging interval adjustable from
      \SIrange{1}{60}{\minute} & D & User survey \\
    \req{req:battery} & Power & Battery life of at least 30 days at a
      \SI{5}{\minute} interval & D & User survey \\
    \req{req:cost} & Cost & Component cost below £40 per unit & D & Market survey \\
    \req{req:size} & Size & Fits within a $\SI{80}{\milli\metre} \times
      \SI{50}{\milli\metre} \times \SI{25}{\milli\metre}$ enclosure & W & Supervisor \\
    \bottomrule
  \end{tabular}
\end{table}

Requirements marked \enquote{D} (demand) must be met; \enquote{W} (wish)
are desirable.
